\subsection{星代数的表示}

设 A 是星代数。以 \(A_{+}\) 表示 A 的正元组成的闭凸锥（I, p. 115, 定义 6）。A 上的线性泛函 \(\lambda\) 是正的，当且仅当 \(\lambda(\mathrm{A}_{+}) \subset \mathbf{R}_{+}\)（I, p. 118, 定理 2）。

命题 7. --- 设 A 是星代数。A 上的每个正线性泛函 \(\lambda\) 都是连续的。

先证明 \(\lambda\) 在 \(\mathrm{A}_{+}\) 与 A 的单位球的交集上有界。反之，存在序列 \((x_{n})_{n\geqslant 1}\) 属于 \(\mathrm{A}_{+}\)，满足 \(\| x_n\| \leqslant 1\) 且 \(\lambda (x_n)\geqslant n\)，对每个整数 \(n\geqslant 1\) 都成立。通项为 \(n^{-2}x_{n}\) 的级数收敛到 A 中的一个元 \(x\)（参见 TG, IV, p. 33, 例 3）。对每个整数 \(\mathrm{N}\geqslant 1\)，由于

\[
\sum_ {n = \mathrm{N} + 1} ^ {+ \infty} \frac {1}{n ^ {2}} x _ {n} \in \mathrm{A} _ {+},
\]

由（I, p. 116, 命题 14）可得

\[
\sum_ {n = 1} ^ {\mathrm{N}} \frac {1}{n} \leqslant \sum_ {n = 1} ^ {\mathrm{N}} \frac {1}{n ^ {2}} \lambda (x _ {n}) = \lambda (x) - \lambda \left(\sum_ {n = \mathrm{N} + 1} ^ {+ \infty} \frac {1}{n ^ {2}} x _ {n}\right) \leqslant \lambda (x),
\]

这与事实矛盾（TG, IV, p. 33, 例 4）。

由于 A 的单位球中的每个元都是 A 中至多四个范数 \(\leqslant 1\) 的正元的系数绝对值不超过 1 的线性组合（I, p. 96, 引理 2 及 I, p. 117, 公式 (4)），所以 \(\lambda\) 连续。

存在一些对合巴拿赫代数，它们具有不连续的正线性泛函（V, p. 493, 习题 3，a)）。

命题 8. --- 设 A 是星代数，a 是 A 的非零元。存在正线性泛函 \(\lambda \in \mathrm{A}_{+}^{\prime}\)，使得 \(\lambda(a^{*}a) > 0\)。

考虑 A 的埃尔米特元组成的实巴拿赫空间 \(A_{h}\)。集合 \(A_{+}\) 是 \(A_{h}\) 中的闭凸尖锥（I, p. 116, 命题 14）。元 \(a^{*}a\) 是正的（I, p. 118, 定理 2）且非零，故埃尔米特元 \(-a^{*}a\) 不是正的。由 EVT, II, p. 42, 推论 5，存在连续实线性泛函 \(\lambda\in A_{h}^{\prime}\)，使得 \(\lambda(-a^{*}a)<0\) 且 \(\lambda(\mathrm{A}_{+})\subset\mathbf{R}_{+}\)。线性泛函 \(\lambda\) 延拓为 A 上的埃尔米特 C-线性泛函（参见 I, p. 98），该泛函是正的，并具有所要求的性质。

定理 2（盖尔范德--奈马克）。--- 设 A 是星代数。存在希尔伯特空间 E 以及 \(\varphi\) 从 A 到 \(\mathcal{L}(\mathrm{E})\) 的等距对合代数态射。若 A 含单位元，则存在这样的保单位元态射。

对每个 \(b \in \mathrm{A} - \{0\}\)，取 A 上的连续正线性泛函 \(\lambda_b\)，使得 \(\lambda_b(b^*b) > 0\)（命题 8）。由于 A 具有逼近单位元（I, p. 121, 命题 18），存在希尔伯特实现 \((\mathrm{E}_b, x_b, \varphi_b)\) 的 \(\lambda_b\)（命题 5）。若 A 含单位元，可设 \(\varphi_b\) 保单位元（同上）。有 \(\| \varphi_b(b) \|^2 = \lambda_b(b^*b) \neq 0\)。

令 E 为各空间 \(E_{b}\)（其中 \(b\) 属于 \(A - \{0\}\)）的外希尔伯特直和。对每个 \(a \in A\)，以 \(\varphi(a) \in \mathcal{L}(E)\) 表示唯一的连续线性映射，其在每个 \(E_{b}\) 上的限制都与 \(\varphi_{b}(a)\) 一致，其中 \(b \in A - \{0\}\)。映射 \(a \mapsto \varphi(a)\) 是对合代数态射；它是单射，因而是等距的（I, p. 112, 命题 9）；若 A 含单位元，还满足 \(\varphi(1) = 1_{\mathrm{E}}\)。证明完毕。

*因此，以星代数为对象、以对合代数态射为态射的范畴，等价于以希尔伯特空间的自同态代数的闭对合子代数为对象、以对合代数态射为态射的范畴。*

